% ---
% RESUMOS
% ---

% resumo em português
\setlength{\absparsep}{18pt} % ajusta o espaçamento dos parágrafos do resumo
\begin{resumo}

%TODO: Reformular essa parte caso seja foco evidenciar que isso é um módulo aplicado em um sistema já construído para a organização ontológica de documentos judiciais

Este trabalho propõe uma prova de conceito para um sistema de tradução automatizada de documentos jurídicos digitalizados, baseado em inteligência artificial. O sistema utiliza técnicas de Processamento de Linguagem Natural (PLN) para interpretar a complexidade da linguagem forense . À medida que processa um número crescente de documentos, o sistema evolui, tornando-se progressivamente mais adaptado às nuances linguísticas e ao jargão judiciário, o que permite traduções mais precisas e contextualmente adequadas, em contraste com tradutores convencionais que tendem a literalismos.

%TODO: (verificações dessas métricas após imlementação / KPI precisa ser elaborado para quantificar a qualidade e efetividade da tradução, talvez pesquisa com "usuários", depende de qual for a quantidade aceitavel de pessoas para uma amostra boa, depende tambem se o perfil precisaria ser judicial)

O principal problema abordado é a acessibilidade e a disseminação internacional de documentos jurídicos brasileiros, facilitando a tradução de conteúdos altamente técnicos e tornando-os acessíveis para um público mais amplo. Aproximadamente 260 milhões de pessoas falam português no mundo, com uma porcentagem significativa fora do Brasil, tornando a tradução de documentos legais uma ferramenta valiosa para a compreensão do sistema jurídico brasileiro em um contexto global.
A tradução automatizada por meio de um sistema PLN, com grau elevado de compreensão contextual, pode também ser reaproveitada e aplicada em diversos cenários onde a linguagem se distancia do uso cotidiano ou das definições de dicionário. Dessa forma, esse sistema se torna uma aliada crucial na interpretação precisa de diferentes linguagens, promovendo uma comunicação mais eficaz em contextos especializados.
Para a implementação do sistema, ferramentas como Python e bibliotecas de PLN foram utilizadas.

%TODO: Elaborar implementação + citação 260 milhões instituto camões https://www.instituto-camoes.pt/images/pdf_noticias/Dados_sobre_a_língua_portuguesa_de_2022.pdf

\textbf{Palavras-chave}: tradução automática, inteligência artificial, documentos jurídicos, processamento de linguagem natural, acessibilidade, python. \end{resumo}






%Descrição geral da empresa (natureza, mercado, processos, etc.), problema-foco atacado no PFC, o que foi feito, principais resultados atingidos, etc.
%Se o documento for escrito em outra língua que não o Português, então é necessário fazer um Resumo Estendido em Português, ao invés deste resumo enxuto.



% resumo em inglês
\begin{resumo}[Abstract]
 \begin{otherlanguage*}{english}
This paper proposes a proof of concept for an automated translation system for digitized legal documents, based on artificial intelligence. The system employs Natural Language Processing (NLP) techniques to interpret the complexity of legal vocabulary. As it processes an increasing number of documents, the system evolves, progressively adapting to linguistic nuances and legal jargon, which enables more precise and contextually appropriate translations, in contrast to conventional translators that tend to resort to literalism.

The primary issue addressed is the accessibility and international dissemination of Brazilian legal documents, facilitating the translation of highly technical content and making it accessible to a broader audience. Approximately 260 million people speak Portuguese worldwide, with a significant percentage residing outside Brazil, making the translation of legal documents a valuable tool for understanding the Brazilian legal system in a global context. Automated translation through an NLP system, with a high degree of contextual understanding, can be reused and applied in various scenarios where language diverges from everyday use or dictionary definitions. In this way, the system becomes a crucial ally in the precise interpretation of different languages, promoting more effective communication in specialized contexts. 
Python and NLP libraries were utilized for the implementation of the system.
   \vspace{\onelineskip}
 
   \noindent 
   \textbf{Keywords}:automated translation, artificial intelligence, legal documents, natural language processing, accessibility, python.
 \end{otherlanguage*}
\end{resumo}

%% resumo em francês 
%\begin{resumo}[Résumé]
% \begin{otherlanguage*}{french}
%    Il s'agit d'un résumé en français.
% 
%   \textbf{Mots-clés}: latex. abntex. publication de textes.
% \end{otherlanguage*}
%\end{resumo}
%
%% resumo em espanhol
%\begin{resumo}[Resumen]
% \begin{otherlanguage*}{spanish}
%   Este es el resumen en español.
%  
%   \textbf{Palabras clave}: latex. abntex. publicación de textos.
% \end{otherlanguage*}
%\end{resumo}
% ---
